%ricordarsi di compilare più di una volta
\documentclass[a4paper,12pt]{letteracdp}

%\usepackage{leaf}
\usepackage[italian]{babel}
\usepackage[T1]{fontenc}
\usepackage[utf8]{inputenc}
\usepackage{layaureo}
\usepackage{lipsum}
\usepackage{fancyhdr}
\usepackage{booktabs}
\usepackage{graphicx}
\usepackage{eurosym}
\usepackage{hyperref}
\hyphenation{Szambelan}

%\lhead{\includegraphics[height=1.5em, width=3.3em]{../../../StileLatex/texmf/tex/commonstuff/logo/leafstuff/logo/logoNoSfondo} \DeclareGraphicsExtensions{.png}}

\author{Davide Castello}

\address{
	\includegraphics[scale=0.25]{logoNoSfondo} \\
	\vspace{0.2cm}
	Ingegneria del Software \\
	Gruppo \textit{Leaf} \\
}

\signature{
	Il Responsabile di Progetto \par
	\includegraphics[scale=0.4, trim=0 1cm 4cm 0cm]{../PianoDiProgetto/firme_componenti/dcastello_firma} \par
	Davide Castello
	}

\date{16 maggio 2016}

\pagenumbering{gobble}
\begin{document}
	\begin{letter}{
		Alla C. Att.ne del committente: \\
		Prof. Vardanega Tullio \\
		Prof. Cardin Riccardo \\
		Dipartimento di Matematica \\
		Università degli Studi di Padova \\
		Via Trieste, 63 \\
		35121 Padova}
		
		\opening{Egregio Prof.~Vardanega Tullio,}
		Con la presente si comunica che il gruppo \textit{Leaf} intende comunicarLe ufficialmente la volontà di partecipare alla \textbf{Revisione di Qualifica} per la realizzazione del prodotto da Lei commissionato, chiamato:
\begin{center}
	\textbf{CLIPS \\ Communication \& Localization with Indoor Positioning Systems}
\end{center}
proposto dall'azienda Miriade S.p.A., rappresentata dalla sig.ra Agnieszka Szambelan.
Può trovare tutti i dettagli relativi alla progettazione di dettaglio del prodotto richiesto all'interno dei seguenti documenti, forniti in allegato:
%
\begin{itemize}
	\item \textit{“Analisi dei requisiti v5.00”} \\(\texttt{RQ/Documents/Esterni/AnalisiDeiRequisiti\_v5.00.pdf});

	\item \textit{“Piano di progetto v5.00”} \\(\texttt{RQ/Documents/Esterni/PianoDiProgetto\_v5.00.pdf});

	\item \textit{“Piano di qualifica v5.00”} \\(\texttt{RQ/Documents/Esterni/PianoDiQualifica\_v5.00.pdf});
	
	\item \textit{“Definizione di prodotto v2.00”} \\(\texttt{RQ/Documents/Esterni/DefinizioneDiProdotto\_v2.00.pdf});
	
	\item \textit{“Manuale utente v1.00”} \\(\texttt{RQ/Documents/Esterni/ManualeUtente\_v1.00.pdf});
	
	\item \textit{"Manuale sviluppatore v1.00”} \\(\texttt{RQ/Documents/Esterni/ManualeSviluppatore\_v1.00.pdf});
	
	\item \textit{“Glossario v5.00”} \\(\texttt{RQ/Documents/Esterni/Glossario\_v5.00.pdf});
	
	\item \textit{“Norme di progetto v4.00”} \\(\texttt{RQ/Documents/Interni/NormeDiProgetto\_v4.00.pdf});

	\item \textit{Verbali delle riunioni interne}	\\(\texttt{RQ/Documents/Interni/Verbali Interni});
	
	\item \textit{Verbali delle riunioni esterne}	\\(\texttt{RQ/Documents/Esterni/Verbali Esterni});
		
\end{itemize}

Oltre alla documentazione, viene consegnato il codice sorgente relativo al prodotto commissionato completo della descrizione di quanto è stato implementato. Il codice sorgente è reperibile al percorso: (\texttt{RQ/Source})
\\\\I dettagli inerenti il consuntivo di periodo sono trattati in maniera approfondita nel \textit{``Piano di Progetto v5.00''} che consegniamo in allegato.

Di seguito si riporta l'elenco dei componenti del gruppo:

\begin{center}
		\begin{tabular}{l c c}
			\toprule
			\textbf{Nominativo} & \textbf{Matricola} & \textbf{Email} \\
			\midrule
			Andrighetto Cristian & 1069767 & black.andrighetto94@gmail.com \\
			%\midrule
			Bicego Eduard & 1069882 & bicegoeduard@gmail.com  \\
			%\midrule
			Castello Davide	& 1073151 &	 castel.davide@gmail.com\\
			%\midrule
			Conti Oscar Elia & 1071039 & conti.oscarelia@gmail.com \\
			%\midrule
			Tavella Federico & 1069038 & federicotavella.7@gmail.com\\
			%\midrule
			Tombolato Andrea & 1069144 & tombolato.andrea@gmail.com	 \\
			%\midrule
			Zanella Marco & 1074420 & zanna0150@gmail.com \\
			\bottomrule
		\end{tabular}
\end{center}
		
		\closing{Rimango a Sua completa disposizione per qualsiasi chiarimento. \\
		RingraziandoLa anticipatamente per la Sua disponibilità, Le porgo distinti saluti.}
		
	\end{letter}

\end{document}
